\begin{ExerciseBlock}
\ExerciseHeading{07}{Reibungsarbeit und Wegabhängigkeit}{work_and_energy_path_dependent_friction_work}
Die Masse $m$ bewegt sich von Punkt $A$ zum Punkt $B$ und zurück.

\begin{center}
	\CachedTikzFigure{content/resources/chapters/work_and_energy/tikz/path_dependent_friction_work.tex}
\end{center}

\begin{enumerate}
	\item Berechnen Sie die Arbeit, die das homogene Gravitationsfeld der Erde für den Hin- und den Rückweg leistet. \[ W_{Grav}^{ges} = W_{Grav}^{AB} + W_{Grav}^{BA} \]

	\item Berechnen Sie die Arbeit, die die (Coulomb'sche) Reibungskraft für den Hin- und den Rückweg leistet. \[ W_{R}^{ges} = W_{R}^{AB} + W_{R}^{BA} \]

	\item Welchen Unterschied beobachten Sie und was bedeutet das für die Existenz einer potentiellen Energie (Lageenergie) für die Gravitation bzw. die Reibekraft.

\end{enumerate}
\end{ExerciseBlock}
\UseExerciseSolution{work_and_energy_path_dependent_friction_work}
